\documentclass[12pt, a4paper]{article}

% ── Shared preamble (src/lecturenotes.sty) ────────────────────────────────────
%   Packages, colors, boxes, \sol/\soleq and math macros all live in
%   src/*.sty. \input@path points LaTeX two folders up to find them.
\makeatletter
\def\input@path{{../../}}
\makeatother
\usepackage{lecturenotes}

% ── Local TikZ styles (match S3/S4) ───────────────────────────────────────────
\usetikzlibrary{decorations.pathreplacing}
\providecommand{\bp}{\mathbf{p}}
\tikzset{
  neuron/.style={circle, draw=mainblue, fill=boxbg,
      minimum size=0.75cm, thick, font=\small},
  preact/.style={circle, draw=mainblue, fill=white,
      minimum size=0.75cm, thick, font=\small},
  input/.style={circle, draw=gray!60, fill=gray!15,
      minimum size=0.75cm, thick, font=\small},
  wire/.style={gray!70, line width=0.5pt},
  layerlabel/.style={font=\small\itshape, gray},
  opbox/.style={draw=orange!70!black, fill=notebg, thick, rounded corners=4pt,
      minimum height=0.9cm, minimum width=1.5cm, font=\small, align=center},
}

% ── Document ──────────────────────────────────────────────────────────────────
\begin{document}

\begin{center}
  {\LARGE \textbf{Section 3.5: The Multi-Layer Perceptron}} \\[0.3em]
  {\large\itshape From a Toy Network to a General Architecture} \\[0.5em]
  {\large AI-IM Class --- Lecture Notes}%
  \ifstudentmode
    \quad {\normalsize\color{gray}[Student Copy]}%
  \fi
\end{center}

\vspace{1em}
\hrule
\vspace{1.5em}

In Sections 3.3 and 3.4 we computed the forward pass and backpropagation
\emph{by hand} for a tiny network: two inputs, two hidden neurons and one
output. Real networks are built from exactly the same pieces, only more of
them. In this section we step back and look at the network as a whole:
\begin{itemize}[nosep]
  \item how the general architecture is organised and how big it can grow,
  \item how to predict several outputs at once,
  \item which activation functions are used in practice,
  \item how gradient descent is run on large datasets, and
  \item how to change the output layer and loss for classification.
\end{itemize}

% ══════════════════════════════════════════════════════════════════════════════
\section{General Architecture}

% ──────────────────────────────────────────────────────────────────────────────
\subsection{Layers and Nodes}

A \textbf{multi-layer perceptron (MLP)}, also called a
\textbf{fully connected feedforward network}, is a stack of \emph{layers}.
Each layer is a column of \emph{nodes} (neurons):

\begin{figure}[H]
  \centering
  \begin{tikzpicture}[x=3.2cm, y=0.95cm, >=Stealth]
    % ── Nodes ────────────────────────────────────────────────────────────────
    \foreach \i [evaluate=\i as \y using {2.5-\i}] in {1,...,4}
    \node[input] (I\i) at (0,\y) {$x_{\i}$};
    \foreach \i [evaluate=\i as \y using {3.5-\i}] in {1,...,6}
    \node[neuron] (A\i) at (1,\y) {};
    \foreach \i [evaluate=\i as \y using {3.5-\i}] in {1,...,6}
    \node[neuron] (B\i) at (2,\y) {};
    \foreach \i [evaluate=\i as \y using {1.5-\i}] in {1,...,2}
    \node[neuron, draw=green!50!black, fill=greenbg] (O\i) at (3,\y)
    {$\yhat_{\i}$};

    % ── Connections ──────────────────────────────────────────────────────────
    \foreach \i in {1,...,4} \foreach \j in {1,...,6}
    \draw[wire] (I\i) -- (A\j);
    \foreach \i in {1,...,6} \foreach \j in {1,...,6}
    \draw[wire] (A\i) -- (B\j);
    \foreach \i in {1,...,6} \foreach \j in {1,...,2}
    \draw[wire] (B\i) -- (O\j);

    % ── Labels ───────────────────────────────────────────────────────────────
    \node[layerlabel, align=center] at (0,-3.4) {Input layer\\(4 features)};
    \node[layerlabel, align=center] at (1,-3.4) {Hidden layer 1\\(6 nodes)};
    \node[layerlabel, align=center] at (2,-3.4) {Hidden layer 2\\(6 nodes)};
    \node[layerlabel, align=center] at (3,-3.4) {Output layer\\(2 outputs)};

    \node[font=\scriptsize, gray] (lbl) at (-0.45, 2.6) {node};
    \draw[->, gray] (lbl) -- (I1);
  \end{tikzpicture}
  \caption{A multi-layer perceptron with four inputs, two hidden layers of six
    nodes each, and two outputs. Every node is connected to every node in the
    next layer.}
  \label{fig:mlp-general}
\end{figure}

\begin{itemize}
  \item The \textbf{input layer} simply holds the feature values
        $x_1, \dots, x_d$. It has no weights and does no computation.
  \item The \textbf{hidden layers} transform the data step by step.
        They are ``hidden'' because we never observe their values in the
        data; the network invents them (recall the learned features
        $h_1, h_2$ of XOR in Section 3.1).
  \item The \textbf{output layer} produces the prediction(s) $\yhat$.
  \item The number of nodes in a layer is its \textbf{width}; the number of
        layers with weights is the network's \textbf{depth}.
\end{itemize}

Our XOR network from Sections 3.3--3.4 is the special case with
2 inputs, one hidden layer of width 2, and 1 output.

% ──────────────────────────────────────────────────────────────────────────────
\subsection{Inside a Node}

Every hidden or output node does the same two sub-steps that we saw in
Section 3.3: a weighted sum, then an activation.

\begin{figure}[H]
  \centering
  \begin{tikzpicture}[>=Stealth]
    % inputs
    \node (h1) at (0,  1.3) {$h_1$};
    \node (h2) at (0,  0)   {$h_2$};
    \node (h3) at (0, -1.3) {$h_3$};
    \node[layerlabel, align=center] at (0, -2.0) {from previous\\layer};

    % node internals
    \node[preact, minimum size=1.1cm] (S) at (3.4, 0) {$\Sigma$};
    \node[neuron, minimum size=1.1cm] (P) at (5.8, 0) {$\phi$};
    \draw[gray!60, dashed, rounded corners=6pt] (2.6,-1.2) rectangle (6.6,1.9);
    \node[layerlabel] at (4.6, -1.55) {one node};

    \draw[->, thick] (h1) -- node[above, font=\scriptsize, pos=0.45] {$w_1$} (S);
    \draw[->, thick] (h2) -- node[above, font=\scriptsize, pos=0.45] {$w_2$} (S);
    \draw[->, thick] (h3) -- node[below, font=\scriptsize, pos=0.45] {$w_3$} (S);
    \node[font=\scriptsize, gray!80!black] (bias) at (3.4, 1.5) {$+\,b$};
    \draw[->, gray] (bias) -- (S);

    \draw[->, thick, mainblue] (S) -- node[above] {$z$} (P);
    \draw[->, thick] (P) -- (8.2, 0) node[right] {$\phi(z)$};

    % formulas
    \node[font=\small] at (4.6, -2.6)
    {$z = w_1 h_1 + w_2 h_2 + w_3 h_3 + b,
        \qquad \text{output} = \phi(z)$};
  \end{tikzpicture}
  \caption{The computation inside a single node. $\Sigma$ forms the
    pre-activation $z$; the activation function $\phi$ is then applied to it.}
  \label{fig:node}
\end{figure}

We write $\phi$ for a generic activation function. In Section 3.3 we used
$\phi = \sigmoid$; the \emph{Activation Functions} part below introduces
other choices.

\paragraph{The whole layer at once.}
Stacking the nodes of layer $\ell$ gives the same matrix form we used in
Section 3.3. With $\bh^{(0)} = \bx$ (the input),
\begin{resultbox}
  \[
    \bz^{(\ell)} = \bW^{(\ell)}\,\bh^{(\ell-1)} + \bb^{(\ell)},
    \qquad
    \bh^{(\ell)} = \phi\bigl(\bz^{(\ell)}\bigr),
    \qquad \ell = 1, 2, \dots, L.
  \]
\end{resultbox}
If layer $\ell-1$ has $n_{\ell-1}$ nodes and layer $\ell$ has $n_\ell$ nodes,
then $\bW^{(\ell)}$ is an $n_\ell \times n_{\ell-1}$ matrix (row $i$ holds the
weights feeding into node $i$) and $\bb^{(\ell)}$ has $n_\ell$ entries.
The forward pass is just this pair of equations repeated layer after layer.

\begin{notebox}
  \textbf{Regression output.} For a regression problem the target $y$ can be
  any real number, so we do \emph{not} squeeze the final output through a
  sigmoid. The output layer uses the \textbf{identity} (``linear'')
  activation: $\yhat = z^{(L)}$.
\end{notebox}

\paragraph{Counting parameters.}
Each node in layer $\ell$ has one weight per incoming node plus one bias, so
\begin{resultbox}
  \[
    \#\text{parameters in layer } \ell
    = \underbrace{n_{\ell-1} \times n_\ell}_{\text{weights}}
    + \underbrace{n_\ell}_{\text{biases}}
    = (n_{\ell-1} + 1)\, n_\ell .
  \]
\end{resultbox}

% ──────────────────────────────────────────────────────────────────────────────
\subsection{More Nodes, More Flexibility}
\label{sec:flexibility}

To see what extra nodes and layers buy us, we fit a single-input,
single-output regression problem: the blue curve below is a wiggly,
noisy function of $x$, and we fit it with increasingly large networks.

\paragraph{Step 0: No hidden layer.}
With the input connected straight to a linear output node, the network is
\[
  \yhat = w\,x + b,
\]
which is ordinary linear regression with $2$ parameters. A straight line
cannot follow the bumps.

\rimg[0.8]{img/fit_linear.png}{No hidden layer: $\yhat = wx + b$
  (2 parameters).}

\paragraph{Step 1: One hidden node.}
Insert one sigmoid node between input and output:
\[
  \yhat = w^{(2)}\,\sigmoid\!\bigl(w^{(1)} x + b^{(1)}\bigr) + b^{(2)}.
\]
Now the fit is a single ``S'' that has been shifted and stretched.
It follows the overall rise of the data but not the bumps
(4 parameters).

\rimg[0.8]{img/fit_1hidden_1node.png}{One hidden node with sigmoid activation
  (4 parameters).}

\paragraph{Step 2: Ten hidden nodes.}
Widen the hidden layer to ten nodes:

\begin{figure}[H]
  \centering
  \begin{tikzpicture}[>=Stealth]
    \node (x) at (-1.4, 0) {$x$};
    \node[input] (In) at (0, 0) {In};
    \node[neuron] (H1) at (3,  2.4) {1};
    \node[neuron] (H2) at (3,  1.2) {2};
    \node[neuron] (H3) at (3,  0)   {3};
    \node at (3, -1.05) {$\vdots$};
    \node[neuron] (H10) at (3, -2.2) {10};
    \node[neuron, draw=green!50!black, fill=greenbg] (Out) at (6, 0) {Out};
    \node (yh) at (7.4, 0) {$\yhat$};

    \draw[->, thick] (x) -- (In);
    \draw[->, thick] (Out) -- (yh);
    \foreach \h in {H1,H2,H3,H10} {
        \draw[wire, ->] (In) -- (\h);
        \draw[wire, ->] (\h) -- (Out);
      }

    \node[font=\scriptsize, gray!80!black, right=0.15cm of H1]
    {$w^{(1)}_1, b^{(1)}_1$};
    \node[font=\scriptsize, gray!80!black, right=0.15cm of H10]
    {$w^{(1)}_{10}, b^{(1)}_{10}$};
    \node[font=\scriptsize, gray!80!black, above=0.25cm of Out, align=center]
    {$w^{(2)}_1, \dots, w^{(2)}_{10}$\\$b^{(2)}$};

    \node[layerlabel] at (3, -3.2) {Hidden layer ($\sigmoid$)};
  \end{tikzpicture}
  \caption{One hidden layer with ten sigmoid nodes and a linear output.}
\end{figure}

The prediction is a \emph{weighted sum of ten S-curves}:
\[
  \yhat = \sum_{j=1}^{10} w^{(2)}_j\,
  \sigmoid\!\bigl(w^{(1)}_j x + b^{(1)}_j\bigr) + b^{(2)}.
\]
Each S-curve can sit at a different place and have a different height and
steepness, so together they can trace the bumps of the data. Count the
parameters:
\[
  \#\text{params}
  = \underbrace{(1+1)\times 10}_{\text{hidden}}
  + \underbrace{(10+1)\times 1}_{\text{output}}
  = \sol[2.5cm]{20 + 11 = 31.}
\]

\rimg[0.8]{img/fit_1hidden_10nodes.png}{One hidden layer, ten sigmoid nodes.
  The fit now follows every bump.}

\paragraph{Step 3: Two hidden layers of 64 nodes.}
Now stack two wide layers: $1 \to 64 \to 64 \to 1$.

\begin{center}
  \renewcommand{\arraystretch}{1.3}
  \begin{tabular}{lccc}
    \toprule
    Layer                     & $n_{\ell-1}$          & $n_\ell$ & Parameters $(n_{\ell-1}+1)\,n_\ell$ \\
    \midrule
    Hidden 1                  & 1                     & 64       & $\sol[2.2cm]{128}$                  \\
    Hidden 2                  & 64                    & 64       & $\sol[2.2cm]{4\,160}$               \\
    Output                    & 64                    & 1        & $\sol[2.2cm]{65}$                   \\
    \midrule
    \multicolumn{3}{l}{Total} & $\sol[2.2cm]{4\,353}$                                                  \\
    \bottomrule
  \end{tabular}
\end{center}

\rimg[0.8]{img/fit_2hidden_64nodes.png}{Two hidden layers of 64 nodes each.}

% ──────────────────────────────────────────────────────────────────────────────
\subsection{Deep Neural Networks}

A network with many hidden layers is called a \textbf{deep neural network};
this is where the name \emph{deep learning} comes from. A common design makes
the layers wide near the input and gradually narrower towards the output:

\begin{figure}[H]
  \centering
  \begin{tikzpicture}[>=Stealth,
      block/.style={draw=mainblue, fill=boxbg, thick, rounded corners=3pt,
          minimum width=0.85cm, font=\small}]
    \node (x) at (-1.3, 0) {$x$};
    \node[input, minimum size=0.9cm] (In) at (0, 0) {In};
    \node[block, minimum height=5.0cm]   (L1) at (1.5, 0) {256};
    \node[block, minimum height=5.0cm]   (L2) at (2.8, 0) {256};
    \node[block, minimum height=2.5cm]   (L3) at (4.1, 0) {128};
    \node[block, minimum height=1.25cm]  (L4) at (5.4, 0) {64};
    \node[block, minimum height=0.75cm]  (L5) at (6.7, 0) {32};
    \node[block, minimum height=0.6cm]   (L6) at (8.0, 0) {16};
    \node[neuron, minimum size=0.9cm, draw=green!50!black, fill=greenbg]
    (Out) at (9.5, 0) {Out};
    \node (yh) at (10.8, 0) {$\yhat$};

    \draw[->, thick] (x) -- (In);
    \draw[thick, gray] (In) -- (L1) (L1) -- (L2) (L2) -- (L3) (L3) -- (L4)
    (L4) -- (L5) (L5) -- (L6) (L6) -- (Out);
    \draw[->, thick] (Out) -- (yh);

    \draw[decorate, decoration={brace, amplitude=6pt, mirror}, gray]
    (1.0,-2.8) -- (8.5,-2.8) node[midway, below=8pt, layerlabel]
    {six hidden layers (block height $\propto$ width)};
  \end{tikzpicture}
  \caption{A deep network: $1 \to 256 \to 256 \to 128 \to 64 \to 32 \to 16
      \to 1$. Each block is a fully connected layer; the number is its width.}
  \label{fig:deep}
\end{figure}

\begin{center}
  \renewcommand{\arraystretch}{1.25}
  \begin{tabular}{lrrr}
    \toprule
    Layer                     & $n_{\ell-1}$            & $n_\ell$ & Parameters            \\
    \midrule
    Hidden 1                  & 1                       & 256      & $512$                 \\
    Hidden 2                  & 256                     & 256      & $65\,792$             \\
    Hidden 3                  & 256                     & 128      & $32\,896$             \\
    Hidden 4                  & 128                     & 64       & $\sol[1.8cm]{8\,256}$ \\
    Hidden 5                  & 64                      & 32       & $\sol[1.8cm]{2\,080}$ \\
    Hidden 6                  & 32                      & 16       & $\sol[1.8cm]{528}$    \\
    Output                    & 16                      & 1        & $\sol[1.8cm]{17}$     \\
    \midrule
    \multicolumn{3}{l}{Total} & $\sol[1.8cm]{110\,081}$                                    \\
    \bottomrule
  \end{tabular}
\end{center}

\rimg[0.8]{img/fit_deep.png}{The deep network of Figure~\ref{fig:deep}.
  The fit is made of short straight segments, the
  signature of the ReLU activation (see Activation Functions below).}

\begin{warnbox}
  \textbf{More parameters is not automatically better.}
  The 31-parameter network already fits this data well. The extra
  parameters of the larger networks mostly go into following the
  \emph{noise}, and they cost more data and computation to train.
  We will return to this trade-off (overfitting) later in the course.
\end{warnbox}

\begin{stepbox}{Exercise 1 \quad Count the parameters}
  A network takes $3$ input features, has two hidden layers of widths
  $5$ and $4$, and produces $2$ outputs. How many weights and biases does
  it have?
  \[
    \#\text{params}
    = \sol[10cm]{(3+1)\cdot 5 + (5+1)\cdot 4 + (4+1)\cdot 2 = 20 + 24 + 10 = 54.}
  \]
\end{stepbox}

% ══════════════════════════════════════════════════════════════════════════════
\section{Multiple Outputs (Regression)}

Nothing forces the output layer to have a single node. If we want to predict
$K$ quantities at once, for example the demand for three products next week,
we simply put $K$ nodes in the output layer, one per target:

\begin{figure}[H]
  \centering
  \begin{tikzpicture}[x=3cm, y=1cm, >=Stealth]
    \foreach \i [evaluate=\i as \y using {2-\i}] in {1,...,3} {
        \node (xin\i) at (-0.55,\y) {$x_{\i}$};
        \node[input] (I\i) at (0,\y) {};
        \draw[->, thick] (xin\i) -- (I\i);
      }
    \node[neuron] (A1) at (1, 1.8) {};
    \node[neuron] (A2) at (1, 0.6) {};
    \node[neuron] (A3) at (1, -0.6) {};
    \node at (1, -1.45) {$\vdots$};
    \node[neuron] (A4) at (1, -2.2) {};
    \foreach \i [evaluate=\i as \y using {2-\i}] in {1,...,3} {
        \node[neuron, draw=green!50!black, fill=greenbg] (O\i) at (2,\y) {};
        \node (yout\i) at (2.55,\y) {$\yhat_{\i}$};
        \draw[->, thick] (O\i) -- (yout\i);
      }
    \foreach \i in {1,...,3} \foreach \j in {1,...,4}
    \draw[wire] (I\i) -- (A\j);
    \foreach \i in {1,...,4} \foreach \j in {1,...,3}
    \draw[wire] (A\i) -- (O\j);

    \node[layerlabel] at (0, -3.0) {Inputs};
    \node[layerlabel] at (1, -3.0) {Hidden layer(s)};
    \node[layerlabel] at (2, -3.0) {Outputs (linear)};
  \end{tikzpicture}
  \caption{A network with three inputs and three outputs. Each output node is
    a linear node with its own weights and bias.}
  \label{fig:multi-output}
\end{figure}

The hidden layers are \emph{shared}: all three predictions are built from the
same learned features. Only the last layer has separate weights for each
output. The loss simply adds up the squared errors over all outputs and all
samples:
\begin{resultbox}
  \[
    J = \frac{1}{n} \sum_{i=1}^{n} \sum_{k=1}^{K}
    \bigl(y^{(i)}_k - \yhat^{(i)}_k\bigr)^2 .
  \]
\end{resultbox}
Backpropagation works exactly as before: each output node receives its own
error signal $\yhat_k - y_k$, and the hidden nodes collect the signals from
\emph{all} outputs they feed into.

\begin{notebox}
  \textbf{Mind the scales.} If one target is measured in thousands and
  another in fractions, the large one dominates the sum. It is good practice
  to standardise each target (and each input) before training.
\end{notebox}

% ══════════════════════════════════════════════════════════════════════════════
\section{Activation Functions}
\label{sec:activation}

The activation $\phi$ is what makes a network more than linear regression:
without it, stacking layers would just multiply matrices together and the
whole network would collapse to a single linear function. Four common choices
are shown below.

% helper: one small axes panel ------------------------------------------------
\newcommand{\actaxes}[4]{% xmin, xmax, ymin, ymax
  \draw[gray!40, ->] (#1,0) -- (#2+0.3,0) node[right, font=\scriptsize, gray] {$z$};
  \draw[gray!40, ->] (0,#3) -- (0,#4+0.3);
}

\begin{figure}[H]
  \centering
  % ── Sigmoid ────────────────────────────────────────────────────────────────
  \begin{minipage}[b]{0.48\textwidth}
    \centering
    \begin{tikzpicture}[x=0.55cm, y=2.6cm]
      \actaxes{-6}{6}{-0.05}{1.05}
      \draw[gray!50, dashed] (-6,1) -- (6,1);
      \foreach \t in {-6,-3,3,6}
      \draw[gray!40] (\t,0) -- (\t,-0.03) node[below, font=\tiny, gray] {$\t$};
      \node[left, font=\tiny, gray, fill=white, inner sep=1pt] at (0,1) {$1$};
      \node[left, font=\tiny, gray, fill=white, inner sep=1pt] at (0,0.5) {$0.5$};
      \draw[mainblue, very thick, smooth, domain=-6:6, samples=80]
      plot (\x, {1/(1+exp(-\x))});
    \end{tikzpicture}\\[2pt]
    \textbf{Sigmoid} \quad $\sigmoid(z) = \dfrac{1}{1+e^{-z}}$\\
    {\small range $(0, 1)$}
  \end{minipage}\hfill
  % ── Tanh ───────────────────────────────────────────────────────────────────
  \begin{minipage}[b]{0.48\textwidth}
    \centering
    \begin{tikzpicture}[x=0.55cm, y=1.3cm]
      \actaxes{-6}{6}{-1.1}{1.1}
      \draw[gray!50, dashed] (-6,1) -- (6,1);
      \draw[gray!50, dashed] (-6,-1) -- (6,-1);
      \foreach \t in {-6,-3,3,6}
      \draw[gray!40] (\t,0) -- (\t,-0.06) node[below, font=\tiny, gray] {$\t$};
      \node[left, font=\tiny, gray, fill=white, inner sep=1pt] at (0,1) {$1$};
      \node[left, font=\tiny, gray, fill=white, inner sep=1pt] at (0,-1) {$-1$};
      \draw[mainblue, very thick, smooth, domain=-6:6, samples=80]
      plot (\x, {tanh(\x)});
    \end{tikzpicture}\\[2pt]
    \textbf{Tanh} \quad $\tanh(z) = \dfrac{e^{z}-e^{-z}}{e^{z}+e^{-z}}$\\
    {\small range $(-1, 1)$}
  \end{minipage}

  \vspace{1.2em}

  % ── ReLU ───────────────────────────────────────────────────────────────────
  \begin{minipage}[b]{0.48\textwidth}
    \centering
    \begin{tikzpicture}[x=0.55cm, y=0.5cm]
      \actaxes{-6}{6}{-1}{6}
      \foreach \t in {-6,-3,3,6}
      \draw[gray!40] (\t,0) -- (\t,-0.15) node[below, font=\tiny, gray] {$\t$};
      \foreach \t in {3,6}
      \node[left, font=\tiny, gray, fill=white, inner sep=1pt] at (0,\t) {$\t$};
      \draw[mainblue, very thick] (-6,0) -- (0,0) -- (6,6);
    \end{tikzpicture}\\[2pt]
    \textbf{ReLU} \quad $\relu(z) = \max(0, z)$\\
    {\small range $[0, \infty)$}
  \end{minipage}\hfill
  % ── Leaky ReLU ─────────────────────────────────────────────────────────────
  \begin{minipage}[b]{0.48\textwidth}
    \centering
    \begin{tikzpicture}[x=0.55cm, y=0.5cm]
      \actaxes{-6}{6}{-1}{6}
      \foreach \t in {-6,-3,3,6}
      \draw[gray!40] (\t,0) -- (\t,-0.15) node[below, font=\tiny, gray] {$\t$};
      \foreach \t in {3,6}
      \node[left, font=\tiny, gray, fill=white, inner sep=1pt] at (0,\t) {$\t$};
      \draw[mainblue, very thick] (-6,-0.6) -- (0,0) -- (6,6);
    \end{tikzpicture}\\[2pt]
    \textbf{Leaky ReLU} \quad $\phi(z) = \max(\alpha z, z)$\\
    {\small range $(-\infty, \infty)$, shown with $\alpha = 0.1$}
  \end{minipage}
  \caption{Four common activation functions.}
  \label{fig:activations}
\end{figure}

\paragraph{Why not always use sigmoid?}
Backpropagation (Section 3.4) multiplies one factor $\phi'(z)$ per layer on
the way back. For the sigmoid, $\sigmoid'(z) = \sigmoid(z)\bigl(1-\sigmoid(z)\bigr)
  \le 0.25$, and it is almost $0$ whenever $|z|$ is large (the flat parts of the
``S''). Through many layers these small factors multiply together, and the
gradient reaching the first layers becomes tiny. This is the
\textbf{vanishing gradient} problem: the early layers barely learn.

\paragraph{ReLU.}
The \textbf{rectified linear unit} has derivative exactly $1$ for $z > 0$, so
gradients pass through active nodes unchanged. It is also very cheap to
compute. ReLU is the default choice for hidden layers in modern networks.
Its weakness: for $z < 0$ the derivative is $0$, so a node whose
pre-activation is always negative stops learning (a ``dead'' node).
\textbf{Leaky ReLU} fixes this by giving the negative side a small slope
$\alpha$ (typically $0.01$--$0.1$).

\begin{center}
  \renewcommand{\arraystretch}{1.3}
  \begin{tabular}{llll}
    \toprule
    Activation & Derivative $\phi'(z)$                  & Range              & Typical use                      \\
    \midrule
    Sigmoid    & $\sigmoid(z)\bigl(1-\sigmoid(z)\bigr)$ & $(0,1)$            & output for binary classification \\
    Tanh       & $1 - \tanh^2(z)$                       & $(-1,1)$           & hidden layers (older networks)   \\
    ReLU       & $1$ if $z>0$, else $0$                 & $[0,\infty)$       & hidden layers (default)          \\
    Leaky ReLU & $1$ if $z>0$, else $\alpha$            & $(-\infty,\infty)$ & hidden layers                    \\
    Identity   & $1$                                    & $(-\infty,\infty)$ & output for regression            \\
    \bottomrule
  \end{tabular}
\end{center}

\begin{notebox}
  \textbf{Look back at the fits.} With sigmoid hidden nodes
  (the fits for Steps 1--3 above) the fitted curve is a sum of
  smooth S-shapes. With ReLU nodes it is a sum of ``hinges'' $\max(0, wx+b)$,
  which is why the deep network's fit is made of straight pieces joined at
  kinks.
\end{notebox}

% ══════════════════════════════════════════════════════════════════════════════
\section{Optimizer: Gradient Descent on Large Datasets}

In Section 3.4 we updated the weights using the gradient from a
\emph{single} XOR input. In general the loss is an average over all $n$
training samples,
\[
  J(\bth) = \frac{1}{n} \sum_{i=1}^{n} J^{(i)}(\bth),
\]
where $\bth$ collects every weight and bias. Its gradient is therefore an
average of $n$ per-sample gradients. When $n$ is in the millions, computing
the exact gradient before \emph{every} small step is far too slow.

\paragraph{Mini-batches.}
The fix is to \textbf{estimate} the gradient from a small random subset of
the data, called a \textbf{mini-batch} of size $B$:
\[
  \nabla J(\bth) \;\approx\; \frac{1}{B} \sum_{i \in \text{batch}}
  \nabla J^{(i)}(\bth),
  \qquad
  \bth \leftarrow \bth - \eta\, \nabla J(\bth).
\]
The estimate is noisy, but it is cheap, so we can take many more steps in the
same amount of time. On average the steps still point downhill.

\begin{figure}[H]
  \centering
  \begin{tikzpicture}[
      sample/.style={draw=gray!70, minimum size=0.55cm, inner sep=0pt},
      >=Stealth]
    \foreach \i in {0,...,11} {
        \pgfmathtruncatemacro{\b}{int(\i/4)}
        \ifcase\b
          \node[sample, fill=boxbg]        at (\i*0.6, 0) {};
        \or
          \node[sample, fill=greenbg]      at (\i*0.6, 0) {};
        \or
          \node[sample, fill=notebg]       at (\i*0.6, 0) {};
        \fi
      }
    \foreach \b/\c in {0/Batch 1, 1/Batch 2, 2/Batch 3} {
        \draw[decorate, decoration={brace, amplitude=5pt, mirror}, gray]
        (\b*2.4-0.27, -0.45) -- (\b*2.4+1.8+0.27, -0.45)
        node[midway, below=6pt, font=\small] {\c};
        \draw[->, gray] (\b*2.4+0.9, -1.35) -- (\b*2.4+0.9, -1.85)
        node[below, font=\scriptsize, black] {update $\bth$};
      }
    \node[font=\small, anchor=west] at (7.4, 0) {$n = 12$ samples (shuffled)};
    \node[font=\small, anchor=west] at (7.4, -0.9) {$B = 4$};
    \node[font=\small, anchor=west] at (7.4, -1.8) {$\Rightarrow$ 3 updates per epoch};
  \end{tikzpicture}
  \caption{One epoch of mini-batch gradient descent. The shuffled training set
    is cut into batches; each batch gives one gradient estimate and one
    update.}
\end{figure}

One full pass through the training data is called an \textbf{epoch}.
Training usually runs for many epochs, reshuffling the data each time.
Depending on the batch size we get three flavours:

\begin{center}
  \renewcommand{\arraystretch}{1.3}
  \begin{tabular}{llll}
    \toprule
    Method              & Batch size  & Updates / epoch     & Character                     \\
    \midrule
    Batch GD            & $B = n$     & $1$                 & exact gradient, slow per step \\
    Stochastic GD (SGD) & $B = 1$     & $n$                 & very noisy, cheap per step    \\
    Mini-batch GD       & $1 < B < n$ & $\lceil n/B \rceil$ & the practical compromise      \\
    \bottomrule
  \end{tabular}
\end{center}

In practice ``SGD'' usually refers to the mini-batch version, with batch
sizes such as $32$, $64$ or $128$.

\begin{stepbox}{Exercise 2 \quad Updates per epoch}
  A training set has $n = 10\,000$ samples and we use $B = 64$.
  How many weight updates happen in one epoch, and in $20$ epochs?
  \[
    \sol[11cm]{\lceil 10\,000 / 64 \rceil = \lceil 156.25 \rceil = 157
      \text{ updates per epoch;}\quad 157 \times 20 = 3\,140 \text{ updates.}}
  \]
\end{stepbox}

\begin{notebox}
  \textbf{Beyond plain SGD.} Plain (mini-batch) SGD uses the same learning
  rate $\eta$ for every weight. \textbf{Adam} (adaptive moment estimation)
  is SGD with an \emph{adaptive learning rate}: it keeps a running average of
  past gradients and adjusts the step size for each weight automatically.
  Adam is the default optimizer in most deep-learning work, and it is the
  topic of the next section.
\end{notebox}

% ══════════════════════════════════════════════════════════════════════════════
\section{Classification}

So far the outputs were real numbers. For \textbf{multiclass classification}
the target is one of $C$ categories, for example a product graded
\emph{A}, \emph{B} or \emph{C}. Three changes turn the regression network into
a classifier: the label encoding, the output layer, and the loss.

% ──────────────────────────────────────────────────────────────────────────────
\subsection{One-Hot Encoding}

We cannot feed the labels ``A, B, C'' into a loss function, and coding them as
$1, 2, 3$ would wrongly suggest that C is ``more'' than A. Instead each label
becomes a vector of length $C$ with a $1$ in the position of the true class
and $0$ elsewhere:

\begin{center}
  \renewcommand{\arraystretch}{1.25}
  \begin{tabular}{ccc}
    \toprule
    Label & Class index & One-hot vector $\by$ \\
    \midrule
    A     & 1           & $[1, 0, 0]$          \\
    B     & 2           & $[0, 1, 0]$          \\
    C     & 3           & $[0, 0, 1]$          \\
    \bottomrule
  \end{tabular}
\end{center}

The output layer gets \textbf{one node per class}, so $C$ output nodes.

% ──────────────────────────────────────────────────────────────────────────────
\subsection{Softmax: From Scores to Probabilities}

The output nodes are linear, so they produce arbitrary real numbers
$z_1, \dots, z_C$, called \textbf{logits} (raw scores). We want instead a
probability for each class. The \textbf{softmax} function does this:
\begin{resultbox}
  \[
    p_k = \operatorname{softmax}(z_1, \dots, z_C)_k
    = \frac{e^{z_k}}{\sum_{j=1}^{C} e^{z_j}},
    \qquad k = 1, \dots, C.
  \]
\end{resultbox}
Exponentiating makes every score positive, and dividing by the sum makes the
scores add up to $1$. So $0 < p_k < 1$ and $\sum_k p_k = 1$. The largest
logit always gets the largest probability, and the predicted class is the one
with the highest $p_k$.

\begin{stepbox}{Exercise 3 \quad Compute a softmax}
  A network outputs the logits $\bz = [2.0,\; 1.0,\; 0.1]$. Compute the class
  probabilities (to three decimals).
  \[
    e^{2.0} = \sol[1.5cm]{7.389}, \quad
    e^{1.0} = \sol[1.5cm]{2.718}, \quad
    e^{0.1} = \sol[1.5cm]{1.105}, \quad
    \text{sum} = \sol[1.5cm]{11.212}
  \]
  \[
    \bp = \bigl[\,\sol[1.3cm]{0.659},\; \sol[1.3cm]{0.242},\;
      \sol[1.3cm]{0.099}\,\bigr]
  \]
\end{stepbox}

% ──────────────────────────────────────────────────────────────────────────────
\subsection{Categorical Cross-Entropy Loss}

In Section 1.2 we derived the binary cross-entropy (BCE) loss for logistic
regression. Its multiclass version is the
\textbf{categorical cross-entropy (CCE)}:
\begin{resultbox}
  \[
    J = -\frac{1}{n} \sum_{i=1}^{n} \sum_{k=1}^{C}
    y^{(i)}_k \,\log p^{(i)}_k .
  \]
\end{resultbox}
Because $\by^{(i)}$ is one-hot, only one term of the inner sum is non-zero:
the one for the true class. The loss is therefore simply
\[
  J = -\frac{1}{n} \sum_{i=1}^{n} \log\bigl(
  \text{probability the model gives to the \emph{true} class of sample } i
  \bigr).
\]
If the model is confident and right ($p \to 1$), the term $-\log p \to 0$.
If it is confident and wrong ($p \to 0$), $-\log p \to \infty$. As in
Section 1.2, $\log$ is the natural logarithm.

\begin{example}
  Two samples with three classes:
  \[
    \begin{array}{lll}
      \by^{(1)} = [0, 1, 0], & \bp^{(1)} = [0.05,\; \mathbf{0.95},\; 0.00]
                             & \text{(true class 2)}                       \\[2pt]
      \by^{(2)} = [0, 0, 1], & \bp^{(2)} = [0.10,\; 0.80,\; \mathbf{0.10}]
                             & \text{(true class 3)}
    \end{array}
  \]
  Picking out the probability of the true class in each row,
  \[
    J = -\frac{1}{2}\bigl(\log 0.95 + \log 0.10\bigr)
    = -\frac{1}{2}\bigl(-0.051 - 2.303\bigr) = 1.177 .
  \]
  Almost all of the loss comes from sample 2, which the model gets wrong.
\end{example}

\begin{stepbox}{Exercise 4 \quad Loss of a single prediction}
  For the softmax output of Exercise 3, compute the CCE loss of this one
  sample if the true class is (a) class 1, (b) class 3.
  \[
    \text{(a)}\; -\log \sol[1.3cm]{0.659} = \sol[1.5cm]{0.417},
    \qquad
    \text{(b)}\; -\log \sol[1.3cm]{0.099} = \sol[1.5cm]{2.313}
  \]
\end{stepbox}

% ──────────────────────────────────────────────────────────────────────────────
\subsection{The Complete Classifier}

Putting it together, the forward pass of a classifier ends with a softmax
and the loss compares its probabilities with the one-hot label:

\begin{figure}[H]
  \centering
  \begin{tikzpicture}[>=Stealth, y=1.1cm]
    % inputs
    \foreach \i [evaluate=\i as \y using {2-\i}] in {1,...,3}
    \node[input] (I\i) at (0,\y) {$x_{\i}$};
    % hidden block
    \node[draw=mainblue, fill=boxbg, thick, rounded corners=4pt,
      minimum width=1.6cm, minimum height=3.4cm, align=center, font=\small]
    (H) at (2.3, 0) {hidden\\layers};
    % logits
    \foreach \i [evaluate=\i as \y using {2-\i}] in {1,...,3}
    \node[preact] (Z\i) at (4.6,\y) {$z_{\i}$};
    \foreach \i in {1,...,3} {
        \draw[wire, ->] (I\i) -- (H.west |- I\i);
        \draw[wire, ->] (H.east |- Z\i) -- (Z\i);
      }
    % softmax
    \node[opbox] (SM) at (6.6, 0) {softmax};
    \foreach \i in {1,...,3} \draw[->, gray] (Z\i) -- (SM.west);
    % probabilities
    \foreach \i [evaluate=\i as \y using {2-\i}] in {1,...,3} {
        \node (P\i) at (8.5,\y) {$p_{\i}$};
        \draw[->, gray] (SM.east) -- (P\i);
      }
    % loss
    \node[opbox] (L) at (10.4, 0) {CCE};
    \foreach \i in {1,...,3} \draw[->, gray] (P\i) -- (L.west);
    \node (y) at (10.4, -2.4) {$\by$ (one-hot)};
    \draw[->, gray] (y) -- (L);
    \draw[->, thick] (L) -- (11.9, 0) node[right] {$J$};

    % labels
    \node[layerlabel] at (4.6, -1.6) {logits};
    \node[layerlabel] at (8.5, -1.6) {probabilities};
  \end{tikzpicture}
  \caption{A three-class classifier. The output layer produces logits; softmax
    turns them into probabilities; the CCE loss compares them with the
    one-hot label.}
  \label{fig:classifier}
\end{figure}

\begin{notebox}
  \textbf{A very clean gradient.} When softmax is followed by CCE, the error
  signal at output node $k$ turns out to be simply
  \[
    \pd{J}{z_k} = p_k - y_k ,
  \]
  the same ``prediction minus target'' form we saw for regression in
  Section 3.4. Backpropagation then continues through the hidden layers
  exactly as before.
\end{notebox}

\paragraph{Binary classification.}
With only two classes we can use a single output node with a sigmoid,
$\hat{p} = \sigmoid(z)$ = probability of class 1, together with the BCE
loss from Section 1.2. This is exactly the setup of our XOR network. (A
two-node softmax with CCE gives the same model; one sigmoid node is just more
economical.)

% ──────────────────────────────────────────────────────────────────────────────
\subsection{Labels and Outputs in Practice}

Deep-learning libraries differ in the details, but two choices always have to
match what the loss function expects:

\begin{enumerate}
  \item \textbf{How the labels are stored.}
        \begin{itemize}[nosep]
          \item As \emph{one-hot vectors}: $\by = [0, 0, 1]$.
          \item As an \emph{integer class index}: $y = 2$ (counting from 0)
                or $y = 3$ (counting from 1). This saves memory when $C$ is
                large; the loss picks out $p_y$ directly, which is the same
                number as $\sum_k y_k \log p_k$.
        \end{itemize}
  \item \textbf{What the network outputs.}
        \begin{itemize}[nosep]
          \item \emph{Probabilities}: the model applies softmax itself,
                e.g.\ $[0.30, 0.15, 0.55]$.
          \item \emph{Logits}: the model stops at $\bz$,
                e.g.\ $[1.5, 0.8, 2.1]$, and the loss function applies
                softmax internally.
        \end{itemize}
\end{enumerate}

\begin{warnbox}
  \textbf{Most libraries prefer logits.} Computing softmax and then $\log$
  separately can overflow ($e^{z}$ for large $z$) or take $\log 0$. A loss
  function that receives the logits can combine both steps in a numerically
  safe way. The common mistake is to apply softmax in the model \emph{and}
  again inside the loss; check which form your loss function expects.
\end{warnbox}

% ══════════════════════════════════════════════════════════════════════════════
\section{Summary}

\begin{center}
  \renewcommand{\arraystretch}{1.35}
  \begin{tabular}{llll}
    \toprule
    Task                     & Outputs & Output activation & Loss         \\
    \midrule
    Regression (one target)  & 1       & identity          & MSE          \\
    Regression ($K$ targets) & $K$     & identity          & MSE (summed) \\
    Binary classification    & 1       & sigmoid           & BCE          \\
    Multiclass ($C$ classes) & $C$     & softmax           & CCE          \\
    \bottomrule
  \end{tabular}
\end{center}

\begin{itemize}
  \item An MLP is a stack of fully connected layers. Each layer computes
        \[\bh^{(\ell)} = \phi\bigl(\bW^{(\ell)} \bh^{(\ell-1)} + \bb^{(\ell)}\bigr).\]
  \item A layer with $n_{\ell-1}$ inputs and $n_\ell$ nodes has
        $(n_{\ell-1}+1)\,n_\ell$ parameters. Width and depth add flexibility,
        but also cost and the risk of fitting noise.
  \item Hidden layers usually use ReLU (or a variant); sigmoid and tanh suffer
        from vanishing gradients in deep networks.
  \item Training uses mini-batch gradient descent: one update per batch, many
        epochs. Adam (next section) adapts the step size per weight.
  \item For classification: one-hot labels, $C$ output nodes, softmax, and
        categorical cross-entropy.
\end{itemize}

\end{document}
